\documentclass[11pt,a4paper]{article}
\usepackage[T1]{fontenc}
\usepackage{lmodern}
\usepackage[utf8]{inputenc}
\usepackage[margin=1.9cm,top=1.8cm,bottom=1.8cm]{geometry}
\usepackage{graphicx}
\usepackage{amsmath,amssymb}
\usepackage{xcolor}
\usepackage{caption}
\usepackage{titlesec}
\usepackage{microtype}
\usepackage{url}
\captionsetup{font=small,labelfont=bf,skip=4pt}
\titlespacing*{\section}{0pt}{12pt}{4pt}
\titleformat{\section}{\normalfont\large\bfseries}{\thesection}{0.7em}{}
\setlength{\parskip}{3pt}
\setlength{\parindent}{0pt}
\newcommand{\fig}[2]{\begin{center}\includegraphics[width=0.97\textwidth]{#1}\\[2pt]
{\small\captionof{figure}{#2}}\end{center}}
\newcommand{\Snaky}{S}
\newcommand{\Z}{\mathbb{Z}}
\definecolor{deep}{RGB}{23,69,110}
\definecolor{piv}{RGB}{217,95,2}
\definecolor{ok}{RGB}{46,125,79}
\definecolor{brk}{RGB}{165,24,42}

\begin{document}

\begin{center}
{\LARGE\bfseries Snaky in 21 Maker moves}\\[2pt]
{\large figures for the overall strategy}\\[8pt]
{\small Companion to OpenAI, \emph{Snaky in 21 Maker moves}, 25 September 2026.\\
\url{https://github.com/openai/math/blob/main/preprints/Snaky-in-21-Maker-moves-September-25-2026/article.pdf}}
\end{center}

\vspace{-2pt}
\section*{The claim, and the shape of the argument}

Maker plays first on the empty infinite square grid; Maker and Breaker alternately claim one
unclaimed cell, ownership is permanent, and Maker wins by owning all six cells of
\[
\Snaky=\{(0,0),(1,0),(2,0),(3,0),(3,1),(4,1)\}
\]
up to translation, reflection and quarter turn. The paper proves that one fixed Maker policy
does this within $21$ Maker claims against every Breaker play. Here is the spine of the proof,
in the order the figures present it.

\textbf{One notion.} A \emph{card} is a triple $(A,T,h)$ with finite $A\subseteq T\subset\Z^2$ and
$h\ge 1$. It is a promise: \emph{at every position with Maker to move, if Maker owns $A$ and
Breaker owns nothing in $T$, Maker has already won or can win within $h$ further claims.}
Nothing is said about how the position arose; extra Maker cells are free, and Breaker stones
outside $T$ are irrelevant. This indifference to history is what lets cards be composed.

\textbf{One rule.} Given cards $(A_i,T_i,h_i)$ and a pivot $p$, put
\[
A=\Big(\bigcup_i A_i\;\cup\;\bigcap_i T_i\Big)\setminus\{p\},
\qquad
T=\{p\}\cup\bigcup_i T_i,
\qquad
h=1+\max_i h_i .
\]
Maker claims $p$ (if $p$ is already owned, Maker claims any free cell instead --- a real move,
and it is counted). Afterwards Maker owns every $A_i$ and the whole common part $\bigcap_i T_i$.
Breaker's single reply $b$ cannot lie in every $T_i$, so some child keeps its envelope Breaker-free,
and that child delivers the win within $h-1$ further claims. One stone, one surviving child:
that is the entire mechanism.

\textbf{One finite object.} Start from the six one-hole copies of $\Snaky$ (height $1$) and apply
the rule $722$ times, in a fixed order, placing each card by a signed permutation and a
translation. The last card needs nothing: $A_{727}=\varnothing$, $|T_{727}|=251$,
$T_{727}\subseteq\{0,\dots,16\}^2$, $h_{727}=21$. With no hypothesis to verify, the promise
holds on the empty board.

\textbf{One descent.} Playing the strategy means walking down that tree: claim the pivot, let
Breaker block one child, follow a surviving one. Each step costs exactly one Maker claim and
lowers the height by one, so the game ends within $21$ claims --- and, because the choices are
fixed in advance, against every Breaker.

Figure~0 is that paragraph as a picture; the remaining figures take the four parts in turn.

\fig{figures/fig0.pdf}{The strategy in one line. Left: the leaves of the construction are the six
one-hole copies of $\Snaky$. Middle: the only rule used anywhere in the proof --- Maker owns
the children's requirements and their common part, and one Breaker stone leaves a child intact.
Right: the card reached at the top of the construction, and a play of it. The construction is
done offline; the game is the descent.}

\section{The target, and why orientation is free}

Everything is invariant under the eight signed permutation matrices: a target is $t+R(\Snaky)$
for $R$ among them. The certificate uses exactly this freedom to line shapes up --- every
reference in the data carries one of the eight codes $s=0,\dots,7$ and a translation.

\fig{figures/fig1.pdf}{The six cells of $\Snaky$, indexed by the lower-left corner of the cell, and
the eight orientations, written in the paper's encoding. Any translate of any of these eight is a
win, so a maker policy may choose its orientation in advance and may switch between them
during play.}

\section{Cards, and the only step in the proof}

Figure~2 shows the rule twice: abstractly, and on the first two cards that use it. Card~6 is a
fork --- Maker owns a column of five, claims the cell beside its middle, and threatens two
distinct completions, only one of which Breaker can take. Card~7 stacks two placed copies of
card~6 whose envelopes overlap in a common column; the intersection of the two envelopes is
that column, so after the pivot $(0,5)$ is removed the requirement shrinks to four cells while
the height grows to $3$. This is the composition pattern repeated $722$ times.

\fig{figures/fig2.pdf}{(a) The rule: Maker claims $p$; the requirements $A_i$ and the whole
intersection $\bigcap_i T_i$ become Maker's; one Breaker stone $b$ misses some envelope. (b) Card
$6$, drawn from the certificate data: dark cells are required, the orange square is the pivot
$(1,4)$, the green squares are the two completions. Breaker can occupy at most one of them.
(c) Card $7$ is two placed copies of card $6$ (grey dashed) whose envelopes share the column
$x=0$; the pivot $(0,5)$ is claimed and the shared column falls out as the new requirement.}

\section{The whole certificate at once}

The data are $728$ cards: the six bases plus $722$ numbered lines, each one application of the
rule, with $1{,}620$ combination nodes in all. Two trends are visible and both matter. As the
height grows, envelopes grow --- a card that promises a win in many moves must be protected
against Breaker stones arriving in the meantime, so it needs a large Breaker-free region. And
required sets shrink, down to the empty set at the top: the higher a card is, the less it assumes
about what Maker has already built.

\fig{figures/fig3.pdf}{All $728$ cards. (a) Envelope size against height: from the six six-cell
bases up to the $251$-cell envelope of card $727$. (b) Required-set size against height, falling
to zero. (c) How many cards sit at each height: a pyramid whose apex is the single card $727$.}

\section{How the top of the tree closes}

The last line of the certificate is a single combination with $32$ placed children, drawn from
six earlier cards in four or eight placements each. Every child requires exactly the cell
$(8,8)$, and the intersection of all $32$ envelopes is exactly $\{(8,8)\}$. So Maker plays the
pivot, and the common part terminates: the union of the requirements is the pivot, the pivot is
deleted, and $A_{727}=\varnothing$. What is left of the six-way choice is the geometry --- the
common part of the envelopes is peeled away group by group, from $35$ cells down to none.

\fig{figures/fig4.pdf}{The $32$ children of card $727$, grouped by the card they come from.
Light cells: one placed envelope from the group. Dark cells: those still common to every
envelope listed so far --- $35,31,21,12,4$ cells, then none. The envelopes share exactly the
pivot $(8,8)$ (orange), so after that claim the last card requires nothing, on a board of
$251$ cells inside $\{0,\dots,16\}^2$.}

\section{Playing it}

Fix an enumeration of the grid and the printed order of every child list. At a card, claim its
pivot; if Maker already owns it, claim the first free cell instead. Stop the moment a copy is
complete. After Breaker's reply, move to the first child whose placed envelope avoids that
cell. The height drops by one at each step, and the pivot of the card reached is the next claim.
Figure~5 is one such play, against a Breaker that greedily blocks as many child envelopes as
possible at each step. In the plays we simulated the descent always finished early --- six to
eight claims against two hundred randomised breakers and against a greedy one --- so the figure
shows the mechanism rather than the worst case.

\fig{figures/fig5.pdf}{One legal play. Shading: how long each cell stays inside the active envelope
(the darker, the longer). Outlines: the active envelope at each level, $251\to169\to40\to25\to11
\to8\to6$ cells, with height $21\to20\to10\to9\to3\to2\to1$. Numbered dots are Maker's claims,
crosses are Breaker's replies, and the green outline is the copy completed on move $7$.}

\section*{What the figures do not show}

The paper also gives a $25$-move construction with reduced envelopes and an explicit opening
cover, a $35$-move construction with four conditional interfaces, and a separate local result
(four stones in a row win against at most three nearby blockers). None of that is drawn here;
the three constructions share the rule of Figure~2 and differ in what they record. Nor do the
figures replace the argument: they are a reading aid for Definitions~2, Lemma~3 and the
$21$-move certificate.

\section*{How these figures were checked}

The geometry in Figures~2, 4 and 5 is not drawn by hand. The certificate file distributed with
the paper (\texttt{verification/certificate.txt}, SHA-256
\texttt{3fa12d36\ldots e3d04}) was verified with the paper's own checker and then re-parsed
independently: all $728$ cards, their sets and heights were recomputed from the rule, matching
$|T_{727}|=251$, $h_{727}=21$ and the intermediate values of Table~1. Figure~4 uses the
placed children of the last line; Figure~5 comes from simulating the Section~4 policy, with the
invariant of Definition~2 (Breaker avoids the active envelope, Maker owns its requirement)
checked at every step, against $200$ randomised breakers and a greedy one.

\end{document}
